% TOP_PROBLEM: {"id": 3309, "title": "Negative association in uniform forests", "category_id": 3, "category": {"id": 3, "name": "graph_theory", "display_name": "Graph Theory", "description": "Problems involving graphs, networks, and their properties.", "slug": "graph-theory", "order_index": 3, "created_at": "2026-07-31T15:26:25.671Z"}, "status": "open", "problem_number": "OPG-1757", "set_id": 12, "statusQualification": "Distinct edges and uniformity over all forest edge sets are explicit; the stated target is pairwise negative correlation."}
% FIRST_POSTED: 2026-10-01
% ENTRY_KIND: statement_only
% UnsolvedMath ID 3309; source catalogue code OPG-1757
% !TeX program = lualatex
% Definition and sources only; this file is not an LLM solution attempt.
\documentclass[11pt,a4paper]{article}
\usepackage[margin=25mm]{geometry}
\usepackage{fontspec}
\setmainfont{lmroman10-regular.otf}[BoldFont=lmroman10-bold.otf,
 ItalicFont=lmroman10-italic.otf,BoldItalicFont=lmroman10-bolditalic.otf]
\usepackage{amsmath,amssymb,mathtools}
\usepackage{unicode-math}
\setmathfont{latinmodern-math.otf}
\AtBeginDocument{\let\setminus\smallsetminus}
\usepackage{xurl,hyperref}
\hypersetup{colorlinks=true,urlcolor=blue}
\setcounter{secnumdepth}{0}
\setlength{\parindent}{0pt}
\setlength{\parskip}{5pt}
\setlength{\emergencystretch}{3em}
\begin{document}
\section*{Negative association in uniform forests}\label{3309}
{\small UnsolvedMath ID 3309 · OPG-1757 · Graph Theory · OpenGarden}
\subsection{Definitions and mathematical statement}
Let $G=(V,E)$ be a finite simple undirected graph. A forest edge set
is a subset $F\subseteq E$ for which the spanning graph $(V,F)$ has
no undirected cycle. Isolated vertices are allowed. Let
$\mathcal F(G)=\{F\subseteq E:(V,F)\text{ is a forest}\}$ and choose
$F$ uniformly from this finite nonempty collection. Forests of all
sizes, including the empty forest, receive equal probability.

For every pair of distinct edges $e,f\in E$, the conjectured inequality
is
\[
                  \mathbb P(e\in F\mid f\in F)
                              \le\mathbb P(e\in F).
\]
The conditional probability is defined because the singleton forest
$\{f\}$ occurs with positive probability. Equivalently,
\[
            \mathbb P(e,f\in F)\le
                         \mathbb P(e\in F)\mathbb P(f\in F).
\]
If $N=|\mathcal F(G)|$, and $N_e,N_f,N_{ef}$ count forests containing
the indicated edges, the exact counting target is
$N\,N_{ef}\le N_eN_f$.

The requirement that $e$ and $f$ be distinct is essential: an event is
positively correlated with itself unless its probability is zero or
one. Also, this distribution is not the uniform spanning-tree
distribution and is not uniform over isomorphism types. The catalogue
title uses ``negative association'', but its displayed target is
specifically pairwise negative correlation of edge indicators; it
does not additionally demand inequalities for every pair of increasing
events depending on disjoint edge sets.
\subsection{Short English statement}
For a uniformly chosen forest edge set, does knowing that one edge is present never increase the probability that a different edge is present?
\subsection{Sources}
\begin{itemize}
\item[S1] ULAM.AI and the UnsolvedMath Contributors, supplied problems.json, record 3309 (OPG-1757), OpenGarden collection. Catalogue identity and source transcription; CC BY 4.0.
\url{https://huggingface.co/datasets/ulamai/UnsolvedMath}.
\item[S2] Open Problem Garden, Negative association in uniform forests. Original problem and discussion; the mathematical formulation below is expanded from the supplied source record.
\url{http://www.openproblemgarden.org/op/negative_association_in_uniform_forests}.
\item[S3] G. R. Grimmett and S. N. Winkler, Negative association in uniform forests and connected graphs (2003), formulation for distinct edge indicators.
\url{https://arxiv.org/abs/math/0302185}.
\end{itemize}
\end{document}
